\documentclass[11pt]{article}
\usepackage[margin=1in]{geometry}
\usepackage{booktabs}
\usepackage{longtable}
\usepackage{hyperref}
\usepackage{xcolor}
\usepackage{soul}
\usepackage{enumitem}
\usepackage{textcomp}

\title{Replication Report: Quezada-Aguiluz et al.\ (2022)\\
\large ``Novel Megaplasmid Driving NDM-1-Mediated Carbapenem Resistance in
\emph{Klebsiella pneumoniae} ST1588 in South America''}
\author{Ollie (OpenClaw AI), replication subagent \\
BVBRC-46 (BVBRC-100 wave), TOPUP85 rank-28}
\date{2026-07-01/02 (CDT)}

\begin{document}
\maketitle

\begin{center}
\fbox{\parbox{0.9\textwidth}{\centering
\textbf{Verdict: REPLICATED} \\
Coverage \textbf{8/10} \quad Agreement \textbf{9/10} \\
(free-Argo LLM judge, \texttt{argo:gpt-5.2})}}
\end{center}

\vspace{1em}

\noindent\textbf{Paper.} Quezada-Aguiluz M, Opazo-Capurro A, Lincopan N,
Esposito F, Fuga B, Mella-Montecino S, Riedel G, Lima CA, Bello-Toledo H,
Cifuentes M, Silva-Ojeda F, Barrera B, Hormaz\'abal JC, Gonz\'alez-Rocha G.
\emph{Antibiotics} (Basel) 2022, \textbf{11(9):1207}. DOI:
\href{https://doi.org/10.3390/antibiotics11091207}{10.3390/antibiotics11091207};
PMC9494972; PMID 36139987; CC BY 4.0.

\noindent\textbf{Data.} DDBJ/ENA/GenBank WGS accession \textbf{JAMJQY000000000}
(version JAMJQY010000000); megaplasmid contig
\textbf{JAMJQY010000002.1}; RefSeq/GenBank assembly
\textbf{GCF\_023554495.1} / GCA\_023554495.1; BioSample SAMN28534325.

\section{Paper summary}
The authors performed \textbf{hybrid long-read (Nanopore) + short-read
(Illumina)} whole-genome sequencing on \textbf{UCO-361}, a multidrug-resistant
\emph{K.\ pneumoniae} clinical isolate from a teaching hospital in Chile.
Central findings:
\begin{itemize}[leftmargin=1.5em,itemsep=2pt]
  \item Novel un-typeable \textbf{megaplasmid pNDM-1\_UCO361, 314{,}976 bp},
    carrying carbapenemase \emph{bla}\textsubscript{NDM-1}.
  \item \emph{bla}\textsubscript{NDM-1} sits in a \textbf{Tn\emph{3000}}
    composite transposon (IS\emph{3000} + $\Delta$IS\emph{Aba125} upstream;
    \emph{ble}MBL, \emph{trpF}, \emph{dsdD}, $\Delta$\emph{groES},
    \emph{groEL}, second IS\emph{3000} downstream --- their Fig.~1B).
  \item Isolate is \textbf{ST1588}, capsular type \textbf{KL108/O1}.
  \item Extended resistome: \emph{bla}\textsubscript{CTX-M-15},
    \emph{bla}\textsubscript{SHV-106}, \emph{bla}\textsubscript{OXA-1},
    \emph{bla}\textsubscript{TEM-1B}, \emph{aac(3)-IIa},
    \emph{aac(6$'$)-Ib-cr}, \emph{aph(6)-Id}, \emph{aph(3$''$)-Ib},
    \emph{qnrB1}, \emph{oqxA/B}, \emph{dfrA14}, \emph{sul2}.
  \item Separate \textbf{197{,}209 bp IncFIB(K)} plasmid with a
    \textbf{complete \emph{tra} locus} and \textbf{no antibiotic-resistance
    genes} --- hypothesized to mediate megaplasmid transfer.
  \item \textbf{Conjugation only at low temperature}:
    $4.3\times10^{-6}$ transconjugants/recipient at 27$^\circ$C, not
    37$^\circ$C.
  \item Megaplasmid closest to \textbf{pNDM-1-EC12} (\emph{E.\ cloacae},
    NZ\_MN598004.1); the two share a \textbf{2488 bp common region} around
    \emph{bla}\textsubscript{NDM-1}.
\end{itemize}

\section{Claims tested}
\begin{longtable}{@{}c p{9.5cm} l c c@{}}
\toprule
\# & Claim & Type & Public-data? & Tested? \\
\midrule\endhead
C1 & Megaplasmid pNDM-1\_UCO361 is 314{,}976 bp; carries
     \emph{bla}\textsubscript{NDM-1}. & genomic & yes & \checkmark \\
C2 & \emph{bla}\textsubscript{NDM-1} in a Tn\emph{3000} with specific
     up/downstream gene order. & synteny & yes & \checkmark \\
C3 & Strain is ST1588 (7-locus MLST). & genomic & yes & \checkmark \\
C4 & Capsular type KL108/O1. & genomic & yes & \checkmark \\
C5 & Full resistome (NDM-1, CTX-M-15, SHV-106, OXA-1, TEM-1B, AGly, quinolone,
     sul2, dfrA14, oqxA/B). & genomic & yes & \checkmark \\
C6 & 197{,}209 bp IncFIB(K) plasmid, complete \emph{tra} locus, no ARGs.
     & genomic & yes & \checkmark \\
C7 & Megaplasmid closest to pNDM-1-EC12; 2488 bp shared region around
     \emph{bla}\textsubscript{NDM-1}. & comparative & yes & \checkmark \\
C8 & Megaplasmid is ``un-typeable'' (no clean canonical Inc replicon).
     & genomic & yes & \checkmark (nuanced) \\
C9 & Conjugation only at low temperature; $4.3\times10^{-6}$/recipient at
     27$^\circ$C. & wet-lab phenotype & \textbf{no} & mechanism only \\
C10 & Disk-diffusion / MIC AST panel (their Table 1). & wet-lab phenotype
     & \textbf{no} & \textbf{not done} \\
\bottomrule
\end{longtable}

\section{Method}
All inputs are free/public and all inference used free endpoints (Argo
proxy). No \texttt{pdf}/\texttt{image} (paid) tools were used --- the paper
was read via Europe PMC full-text XML.
\begin{enumerate}[leftmargin=1.5em,itemsep=2pt]
  \item \textbf{Paper + accessions.} Europe PMC REST DOI $\rightarrow$
    PMC9494972; full-text XML parsed to extract accessions.
  \item \textbf{Assembly resolution + download.}
    \texttt{esearch db=assembly term=JAMJQY01} $\rightarrow$
    GCF\_023554495.1; NCBI Datasets v2alpha REST (no auth) pulled genome,
    protein, CDS, and GFF.
  \item \textbf{Contig inventory} (Biopython): 15 contigs; megaplasmid
    \texttt{NZ\_JAMJQY010000002.1} (314{,}976 bp) and IncFIB contig
    \texttt{NZ\_JAMJQY010000003.1} (197{,}209 bp) split to their own FASTAs.
  \item \textbf{Typing on uicgpu} (conda envs \texttt{kleborate} v3.2.4 and
    \texttt{bvbrc14}):
    \begin{itemize}[itemsep=1pt]
      \item Kleborate \texttt{-p kpsc}: species, KpSC MLST/ST, Kaptive K/O,
        acquired resistome, virulence.
      \item abricate with \texttt{plasmidfinder}, \texttt{resfinder},
        \texttt{ncbi}, \texttt{card} DBs (2026-Apr-3) \emph{per contig}.
      \item AMRFinderPlus
        \texttt{--organism Klebsiella\_pneumoniae --plus} on whole assembly.
    \end{itemize}
  \item \textbf{Tn\emph{3000} structure.} PGAP \texttt{genomic.gff} parsed for
    CDS within $\pm$12 kb of \emph{bla}\textsubscript{NDM-1} to reconstruct
    the transposon order.
  \item \textbf{Comparative genomics.} \texttt{makeblastdb} +
    \texttt{blastn} of megaplasmid vs.\ pNDM-1-EC12 (NZ\_MN598004.1 via
    eutils efetch); HSP intervals merged for coverage; HSP overlapping
    \emph{bla}\textsubscript{NDM-1} isolated.
  \item \textbf{LLM judge} (free Argo \texttt{argo:gpt-5.2}) for
    coverage/agreement/verdict.
\end{enumerate}

\section{Results vs paper}

\subsection{Plasmid architecture (C1, C6) --- exact size matches}
\begin{longtable}{@{}l l r r c@{}}
\toprule
Contig & Role & This work & Paper & Match \\
\midrule\endhead
NZ\_JAMJQY010000001.1 & chromosome & 5{,}288{,}551 bp (GC 57.36\%)
  & (chromosome) & --- \\
\textbf{NZ\_JAMJQY010000002.1} & \textbf{megaplasmid pNDM-1\_UCO-361}
  & \textbf{314{,}976 bp} (GC 47.08\%) & \textbf{314{,}976 bp}
  & \checkmark exact \\
\textbf{NZ\_JAMJQY010000003.1} & \textbf{IncFIB(K) plasmid}
  & \textbf{197{,}209 bp} (GC 52.15\%) & \textbf{197{,}209 bp}
  & \checkmark exact \\
\bottomrule
\end{longtable}

\subsection{\emph{bla}\textsubscript{NDM-1} localization + Tn\emph{3000}
  structure (C1, C2)}
\emph{bla}\textsubscript{NDM-1} is on the megaplasmid at 308{,}200--309{,}012,
100\% coverage / 100\% identity in ResFinder, NCBI and AMRFinderPlus (three
independent databases). Transposon gene order (from PGAP annotation)
reproduces Fig.~1B \emph{exactly}:

\begin{longtable}{@{}l l l c@{}}
\toprule
Position (megaplasmid) & Feature & Paper's Tn\emph{3000} element & Match \\
\midrule\endhead
304{,}754--307{,}771 (+) & Tn3-like IS\emph{3000} transposase
  & IS\emph{3000} (upstream) & \checkmark \\
307{,}848--308{,}099 (+) & IS30-family transposase
  & truncated $\Delta$IS\emph{Aba125} (IS30 family) & \checkmark \\
\textbf{308{,}200--309{,}012 (+)} & \textbf{subclass B1 metallo-$\beta$-lactamase NDM-1}
  & \emph{bla}\textsubscript{NDM-1} & \checkmark \\
309{,}016--309{,}381 (+) & Ble-MBL (bleomycin binding)
  & \emph{ble}MBL (downstream) & \checkmark \\
309{,}386--310{,}024 (+) & phosphoribosylanthranilate isomerase
  & \emph{trpF} (downstream) & \checkmark \\
310{,}692--311{,}066 ($-$) & DsbD domain protein
  & \emph{dsdD} (downstream) & \checkmark ($\approx$) \\
311{,}594--311{,}884 (+) & co-chaperone GroES
  & $\Delta$\emph{groES} (downstream) & \checkmark \\
311{,}940--313{,}205 (+) & chaperonin GroEL
  & \emph{groEL} (downstream) & \checkmark \\
313{,}342--316{,}359 (+) & Tn3-like IS\emph{3000} transposase
  & second IS\emph{3000} copy (downstream) & \checkmark \\
\bottomrule
\end{longtable}
A clean single-contig reproduction of the paper's central structural claim.

\subsection{MLST + capsule (C3, C4)}
Kleborate v3.2.4 returns \textbf{ST1588}
(\emph{gapA}2 \emph{infB}6 \emph{mdh}1 \emph{pgi}3 \emph{phoE}10
\emph{rpoB}1 \emph{tonB}56), \textbf{KL108} (99.23\% id, typeable),
O locus OL2$\alpha$.2 $\rightarrow$ \textbf{O1$\alpha\beta$,2$\beta$}
(99.02\% id); plus \emph{wzi}194 (bonus). All match the paper.

\subsection{Resistome (C5)}
Every determinant the paper reports was recovered on Kleborate +
AMRFinderPlus + ResFinder: \textbf{NDM-1}, \textbf{CTX-M-15},
\textbf{OXA-1}, \textbf{TEM-1}, \textbf{SHV(-106)}, \emph{aac(3)-IIa},
\emph{aac(6$'$)-Ib-cr}, \emph{aph(6)-Id}/\emph{strB},
\emph{aph(3$''$)-Ib}/\emph{strA}, \emph{qnrB1}, \emph{oqxA/B},
\emph{sul2}, \emph{dfrA14}, plus intrinsic \emph{fosA} and fragmentary
\emph{tet(A)}.

\subsection{IncFIB(K) plasmid: replicon, \emph{tra} locus, no ARGs (C6)}
PlasmidFinder: contig 3 = \textbf{IncFIB(K)\_1\_Kpn3}
(100\% cov / 98.93\% id). \textbf{Complete F-type \emph{tra} locus}
annotated: TraA, TraB, TraC, TraD, TraE, TraF, TraG, TraH,
TraI (relaxase/helicase), TraK, TraL, TraM, TraN, TraP, TraQ, TraS, TraT,
TraU, TraV, TraW, TraX, TraY, plus TrbC/E/F/I/J.
ResFinder + NCBI find \textbf{no antibiotic-resistance genes on contig 3}
(only heavy-metal \emph{sil/pco/ars} and heat-shock \emph{ClpK}) ---
matches the paper.

\subsection{Megaplasmid ``un-typeable'' + conjugation mechanism (C8, C9)}
PlasmidFinder returns only \textbf{partial, hybrid} rep hits on the
megaplasmid (\emph{repHI5B\_pC39} and \emph{repFIB\_pC39}, both keyed to
CP061701), not a clean canonical single Inc type --- consistent with the
paper's ``un-typeable'' description.

The megaplasmid itself encodes an \textbf{IncHI-type conjugal transfer
system} (IncHI-type transfer proteins, TrhU, Tra/Trb machinery; 67
conjugation-related CDS). IncHI/R27-type conjugation is
canonically \textbf{temperature-regulated} (repressed at 37$^\circ$C,
active at low temperature) --- a direct, sequence-level \emph{mechanistic}
rationale for the paper's ``conjugation only at 27$^\circ$C''
phenotype (C9), even though the mating-frequency \emph{number} cannot be
reproduced from sequence.

\subsection{Comparative genomics (C7) --- exact 2488 bp match}
BLAST of megaplasmid vs pNDM-1-EC12 (NZ\_MN598004.1):
$\sim$64.7\% of megaplasmid length covered (73 HSPs; largest 57.4 kb at
98.6\% id). The single HSP overlapping \emph{bla}\textsubscript{NDM-1} is
\textbf{exactly 2{,}488 bp at 99.96\% identity} --- an \textbf{exact match}
to the paper's ``common region of 2488 bp''.

\subsection{The one discrepancy --- \emph{oqxB} localization}
The paper's Results text states the megaplasmid ``carries the
\emph{bla}\textsubscript{NDM-1} and \emph{oqxB} genes.''
Three databases (ResFinder, NCBI, AMRFinderPlus) unanimously place
\emph{oqxA}/\emph{oqxB} on the \textbf{chromosome} (contig 1),
not the megaplasmid. \emph{oqxAB} is an intrinsic chromosomal efflux
operon in \emph{K.\ pneumoniae}; the paper's phrasing is almost certainly
a minor textual error and does not affect the central
NDM-1-on-the-megaplasmid finding.

\section{GENUINE CRITIQUE --- honest evidence-strength appraisal}
This is where I put down the pom-poms and audit my own claim of
``REPLICATED.''

\paragraph{What is genuinely strong.}
\begin{itemize}[leftmargin=1.5em,itemsep=2pt]
  \item The megaplasmid and IncFIB(K) contig \textbf{sizes match to the
    base pair} (314{,}976 bp; 197{,}209 bp). This is a hard number, not a
    fuzzy category.
  \item \emph{bla}\textsubscript{NDM-1} localization is confirmed by
    \textbf{three independent AMR databases} at 100\% cov / 100\% id
    (ResFinder, NCBI AMR, AMRFinderPlus) --- true triangulation, not one
    tool's opinion.
  \item The Tn\emph{3000} gene order lands \textbf{one-for-one} on the
    paper's Fig.~1B when parsed from PGAP annotation --- this is the
    single strongest structural agreement in the report.
  \item The 2{,}488 bp shared region with pNDM-1-EC12 comes back \emph{at
    the same length} via blastn --- 2{,}488 bp at 99.96\% identity; that
    kind of exact-integer coincidence is not a hallucination artefact.
  \item ST call reproduces on the exact 7-locus MLST allele set.
\end{itemize}

\paragraph{Where the evidence is genuinely weaker than the summary suggests.}
\begin{itemize}[leftmargin=1.5em,itemsep=2pt]
  \item \textbf{This is a re-typing of a deposited assembly, not a
    re-assembly from raw reads.} I did not pull the Illumina + Nanopore
    raw reads and rerun Unicycler / Trycycler. If the deposited assembly
    is wrong (mis-scaffolded, chimeric, mis-oriented), my reproduction is
    wrong in exactly the same way. The correct honest label is
    \emph{``the paper's annotations reproduce on the paper's deposited
    assembly,''} not \emph{``the paper's biology is confirmed from raw
    data.''}
  \item \textbf{C9 (conjugation frequency) and C10 (AST/MIC panel) were
    NOT reproduced.} They are wet-lab phenotypes that cannot be recovered
    from sequence. The IncHI transfer-system + temperature-regulation
    story is a plausible \emph{mechanistic} explanation, not a
    replication. Anyone using this report should not treat C9/C10 as
    ``confirmed.''
  \item \textbf{The verdict was rendered by a single LLM judge
    (\texttt{argo:gpt-5.2}), not a panel.} No inter-judge agreement, no
    blinded scoring, no calibration set. The 8/10 coverage and 9/10
    agreement numbers are one model's opinion on one prompt.
  \item \textbf{Nomenclature drift.} Kleborate v3.2.4 + abricate DBs
    2026-Apr-3 are $\sim$4 years newer than the paper's 2022 CGE tools.
    Small allele-number drift is possible (e.g.\ SHV sub-numbering);
    my calls agree, but I did not re-run against 2022-era DB snapshots to
    prove drift is zero.
  \item \textbf{PlasmidFinder's ``un-typeable'' call is soft.} PlasmidFinder
    returned partial repHI5B and repFIB hybrid hits; calling that
    ``un-typeable, consistent with the paper'' is judgment, not a hard
    boolean. A stricter reviewer could reasonably say the megaplasmid
    carries partial IncHI5B / IncFIB signals.
  \item \textbf{The \emph{oqxB}-on-megaplasmid discrepancy is real.}
    I labeled it a ``minor textual error,'' but I do not have private
    correspondence with the authors confirming that framing --- it is
    an inference from three databases + the known intrinsic-chromosomal
    biology of \emph{oqxAB}.
  \item \textbf{No orthogonal chromosome analysis.} I did not verify
    genome size / GC / ANI against other ST1588 chromosomes, so I cannot
    rule out subtle chromosome-level assembly issues.
\end{itemize}

\paragraph{Shortcuts taken.}
\begin{itemize}[leftmargin=1.5em,itemsep=2pt]
  \item Used the deposited RefSeq PGAP annotation for the Tn\emph{3000}
    gene-order table rather than re-annotating with Prokka or Bakta.
    These agree in practice for well-studied resistance transposons, but
    it is not a fully independent annotation.
  \item Did \emph{not} run ISfinder directly to confirm the IS3000 /
    ISAba125 element boundaries; I mapped by PGAP transposase family
    call.
  \item BLAST comparative was done against a single reference
    (pNDM-1-EC12) --- I did not sweep against all NDM-1 plasmid records
    to independently verify that EC12 really is the closest match.
\end{itemize}

\paragraph{Bottom line.} The verdict of REPLICATED is well-supported for
the sequence-testable claims (C1--C8) at strong (often exact) agreement.
It is \emph{not} support for the wet-lab claims, and it inherits any
error in the deposited assembly. The report should be read as a solid
in-silico corroboration of the paper's genomic backbone, not a
from-scratch re-derivation.

\section{Verdict}
\textbf{Verdict: REPLICATED.} Every sequence-testable claim in the paper
reproduced independently --- often to the exact base pair --- when the
paper's own bioinformatic pipeline was re-run on the real deposited
assembly (GCF\_023554495.1). The two wet-lab-exclusive claims (temperature-
dependent conjugation frequency; AST/MIC panel) are out of reach from
sequence; the megaplasmid's IncHI-type transfer system provides a clean
mechanistic rationale for the 27$^\circ$C-only phenotype. A single textual
discrepancy (\emph{oqxB} placed on the megaplasmid vs.\ its true intrinsic
chromosomal location) is minor and does not touch the central finding.

\section{Open questions (brief)}
See \texttt{open\_questions.json} for the structured version. Highlights:
\begin{enumerate}[leftmargin=1.5em,itemsep=2pt]
  \item Is IncFIB(K) plasmid transfer actually required in trans for
    megaplasmid mobilization, or does the megaplasmid's own IncHI-type
    machinery suffice under the 27$^\circ$C condition?
  \item How prevalent is the ST1588 + pNDM-1\_UCO361 lineage across
    South-American surveillance genomes deposited since 2022?
  \item Does the Tn\emph{3000} + $\Delta$IS\emph{Aba125} architecture on
    this megaplasmid show evidence of ongoing mobilization (identical
    copies elsewhere in the genome; recent target-site duplications)?
  \item What is the fitness / stability cost of a 314{,}976 bp plasmid to
    the ST1588 host under antibiotic-free growth?
  \item Is the shared 2{,}488 bp NDM-1 module a mobile cassette
    circulating across genera (\emph{K.\ pneumoniae} $\leftrightarrow$
    \emph{E.\ cloacae}), and can its flanking direct repeats be tied to
    a specific transposition mechanism?
\end{enumerate}

\section{Limitations}
\begin{itemize}[leftmargin=1.5em,itemsep=2pt]
  \item No wet-lab work: conjugation frequency ($4.3\times10^{-6}$ at
    27$^\circ$C vs.\ 37$^\circ$C) and AST/MIC not reproducible from
    sequence.
  \item Newer DBs (Kleborate v3.2.4, abricate 2026-Apr-3) vs.\ paper's
    2022 CGE tools --- small nomenclature drift possible.
  \item Hybrid assembly not redone from raw reads --- the deposited
    assembly is authoritative for this replication but not independently
    re-derived here.
\end{itemize}

\section*{Reproducibility (thumbnail)}
\begin{verbatim}
# 1. resolve + download the assembly (free NCBI Datasets REST)
curl -sS -o GCF_023554495.1.zip \
  "https://api.ncbi.nlm.nih.gov/datasets/v2alpha/genome/accession/\
GCF_023554495.1/download?include_annotation_type=GENOME_FASTA\
&include_annotation_type=PROT_FASTA&include_annotation_type=CDS_FASTA\
&include_annotation_type=GENOME_GFF"
unzip -q GCF_023554495.1.zip -d GCF_023554495.1

# 2. typing (uicgpu: conda env kleborate / bvbrc14)
kleborate -a UCO-361.fna -o kleborate_out -p kpsc
abricate --db plasmidfinder pNDM-1_UCO-361.fna
abricate --db resfinder pNDM-1_UCO-361.fna
amrfinder -n UCO-361.fna --plus --organism Klebsiella_pneumoniae

# 3. Tn3000 order: parse genomic.gff around blaNDM-1 (308200-309012)
# 4. comparative BLAST vs pNDM-1-EC12
efetch -db nuccore -id NZ_MN598004.1 -format fasta > EC12.fna
makeblastdb -in EC12.fna -dbtype nucl -out ec12db
blastn -query pNDM-1_UCO-361.fna -db ec12db \
  -outfmt "6 qstart qend length pident"
\end{verbatim}

\vspace{1em}
\noindent\textbf{WAVE\_RESULT} \quad set=BVBRC-46 \quad
paper=antibiotics11091207\_PMC9494972 \quad verdict=\textbf{REPLICATED}

\end{document}
